\section{Observe the boundaries before comparing work}
Begin with four separated writable words and stop after each quarter-round group.
Next stop the scalar block after its first column round, then its first diagonal
round. For the vector version, inspect each row immediately before selection,
immediately after selection and after the inverse selection. These three views
separate a word-arithmetic error from a lane-order error.

The free-writing contract concerns output, preservation and owned accesses. It does
not require the illustrated scratch registers, local labels or exact instruction
sequence. A working implementation can retain words in registers, choose a different
stack frame or express a rotate through shifts. The register-specific observations
apply to the worked construction, not to every equivalent solution.

Use the same 64-byte initial state when measuring the two block entries. Start the
instruction trace at the function entry and stop at its return continuation. Include
its frame saves and restores; exclude caller setup, assembly and session creation.
The scalar prototype calls its quarter-round eighty times, constructing four pointers
for each call. The MSA prototype keeps rows in registers and selects diagonals with
six shuffles per double-round. Their work difference therefore includes register
residency and call overhead as well as lane arithmetic.

The verified prototypes execute the following work at that boundary. The second
scalar version expresses each word rotate with two shifts and an OR. All three
block versions retain the same input/output contract.

\begin{table}[htbp]\centering
\caption{Executed instructions for one block, including normal return paths.}
\tagpdfsetup{table/header-rows={1}}
\begin{tabular}{lr}\toprule
Implementation & Instructions\\\midrule
Scalar quarter-round calls, rotates & 2816\\
Scalar quarter-round calls, shifts and OR & 3456\\
MSA rows and lane shuffles & 510\\
\bottomrule\end{tabular}\end{table}

The cumulative scalar version retains supplied status checks after all eighty
nested calls. Each check adds a branch and a no-op, accounting for 160 instructions
beyond the earlier prototype. These checks preserve a clear development return while
the program is incomplete; they are included in the measured normal path here.

There are 320 word rotations in eighty quarter-rounds. Replacing each single
rotate with three instructions adds 640 instructions, matching the difference
between the two scalar counts. That comparison isolates the rotate expression;
the MSA comparison changes more of the implementation. The block instruction
counts are independent of the tested input words because the loop bounds and
selection order are fixed.

Count the XOR wrapper separately at message lengths zero, one, 63, 64, 65 and 4096.
The jump from 64 to 65 includes a whole additional block even though only one extra
message byte is transformed. A graph of instructions per byte should identify that
boundary instead of treating the curve as smooth. These are simulated instruction
counts. They establish neither host elapsed-time speedup nor processor cycle counts.

After correctness, try inlining the scalar quarter-round or retaining the scalar
working state in registers. Compare those changes with the original scalar function
before attributing the whole difference to MSA. Keep the same input/output ownership
and counter behavior for each experiment. A shorter loop that reads past a final
fragment has changed the program's contract.

\begin{MARSExercise}{Explain a misplaced lane}
Suppose row one still contains indices \([4,5,6,7]\) during the diagonal round,
while rows two and three are selected correctly. State the four indices used by lane
zero. Identify which intended diagonal input is missing. Explain why restoring the
rows afterward cannot undo the arithmetic already performed on the wrong selection.
\end{MARSExercise}
\begin{MARSExercise}{Design a retained keystream context}
Extend the whole-message interface on paper so a one-byte call followed by a 64-byte
call consumes consecutive keystream bytes. Name the retained block, byte position,
next counter and exhaustion state. Describe how the context represents the last
counter's remaining bytes without allowing another block after they are consumed.
Do not implement it by restarting the supplied state for each call.
\end{MARSExercise}
